\section{Complete physical phase arithmetic}
\label{sec:complete-phase-arithmetic}
The roof and constant parts of a physical collision phase have already
been determined in Theorem~\ref{thm:physical-phase-rigidity}. We now
determine its two spatial parts. The argument uses an atom of the phase
on a positive-measure product set, not a value at a periodic point.
It applies to arbitrary real spatial frequencies, including irrational
ones.

Write
\[
 \Omega=\T^2\times\T\times\R,\qquad
 \phi_{R,\omega}=u\cdot\kappa_R+s+t\tau_R,
 \quad \omega=(u,s,t).
\]
The three torus coordinates have period $2\pi$. The origin of $\Omega$
therefore means $u\in2\pi\Z^2$, $s\in2\pi\Z$, $t=0$.

\begin{lemma}[Separation of a proper integer subgroup]
\label{lem:integer-subgroup-separation}
If $L$ is a proper subgroup of $\Z^d$, then there are a prime $p$ and a
nonzero $a\in(\Z/p\Z)^d$ with $a\cdot L=0$ modulo $p$.
\end{lemma}
\begin{proof}
An integer subgroup has an integer basis. If its rank is less than $d$,
choose a nonzero integer row annihilating it and a prime not dividing
all coordinates of that row. If it has full rank, its index is the
absolute determinant of a basis matrix. This integer exceeds one.
Modulo a prime divisor of the index, the basis matrix is singular and
has a nonzero left null vector. These two constructions prove the
claim without a rationality assumption on the character that might
have defined $L$.
\end{proof}

\begin{theorem}[Complete measurable physical phase rigidity]
\label{thm:complete-physical-phase}
For every $R\in I$, a measurable $q:M\to\mathbb S^1$ satisfies
\begin{equation}\label{eq:complete-physical-phase}
 q\circ T_R=e^{i\phi_{R,\omega}}q
\end{equation}
if and only if $\omega=0$ in $\Omega$ and $q$ is constant almost
everywhere.
\end{theorem}
\begin{proof}
Theorem~\ref{thm:physical-phase-rigidity} gives $t=0$ and
$s\in2\pi\Z$. Both holonomy formulas
\eqref{eq:measurable-stable-holonomy} and
\eqref{eq:measurable-unstable-holonomy} consequently have right side
one. In the product coordinates of
Lemma~\ref{lem:regular-collision-product}, the first says that $q$
is almost everywhere a function only of the unstable transversal
coordinate, and the second that it is almost everywhere a function
only of the stable transversal coordinate. Product measure equivalence
and Fubini imply that $q=c$ almost everywhere on $\mathcal P$, for
one $c\in\mathbb S^1$.

The saturation $\bigcup_{j\ge0}T_R^j\mathcal P$ has full measure by
ergodicity and recurrence. Iterating \eqref{eq:complete-physical-phase}
from $\mathcal P$ therefore shows that
\begin{equation}\label{eq:countable-phase-range}
 q(x)/c\in\{e^{iu\cdot k}:k\in\Z^2\}
       \quad\hbox{for almost every }x.
\end{equation}
Let $L=\{k\in\Z^2:u\cdot k\in2\pi\Z\}$. The map
$[k]\mapsto e^{iu\cdot k}$ identifies the countable group $\Z^2/L$
with the displayed range. Its inverse on that countable range is
measurable. Thus \eqref{eq:countable-phase-range} defines a measurable
$h:M\to\Z^2/L$ with
\begin{equation}\label{eq:countable-phase-cocycle}
 h(T_Rx)=h(x)+[\kappa_R(x)].
\end{equation}
This is a countable-valued lift, not a real logarithm with a chosen
branch. No integrability of the lift is required.

If $L\ne\Z^2$, apply Lemma~\ref{lem:integer-subgroup-separation}.
The resulting nonzero row $a$ defines a homomorphism
$\chi:\Z^2/L\to\Z/p\Z$. On the actual physical billiard modulo
$p\Lambda$ the collision map is
\[
 \widetilde T_R(x,\ell)
 =(T_Rx,\ell+\kappa_R(x)\bmod p),\qquad
 \ell\in(\Z/p\Z)^2.
\]
It preserves $\nu$ times the uniform sheet probability and is ergodic.
Indeed its lift to the plane has exactly the original disk array, so
its free-flight bound is unchanged. Its finitely many disjoint circular
scatterers have positive curvature, and the complement of these disks
in the torus is connected. Thus the finite-horizon billiard theorem
\cite[Section 8.1]{Young} applies. Equivalently use the rectangular
Euclidean cover with periods $(p,0)$ and $(0,p\sqrt3)$, with $2p^2$
scatterers, and then its quotient. No affine deformation or independent
sheet dynamics is introduced.

Equation~\eqref{eq:countable-phase-cocycle} makes
\[
 \exp\!\left(\frac{2\pi i}{p}
             [a\cdot\ell-\chi(h(x))]\right)
\]
invariant under this collision map. The brackets in this formula are
interpreted modulo $p$. The function has modulus one and mean zero
upon averaging over $\ell$, since $a\ne0$. This contradicts ergodicity.
Hence $L=\Z^2$, or $u\in2\pi\Z^2$. The original phase equation now
says that $q$ is invariant, so it is constant. The converse is immediate.
\end{proof}

\subsection{Finite positive-measure return witnesses}
This consequence will replace a nonquantitative appeal to absence of
exact eigenfunctions in the next section. The auxiliary set below is
not the return section $Y_R^*$ of the physical record.

\begin{lemma}[A finite generating family of product-set returns]
\label{lem:finite-return-witnesses}
Fix $R$ and a positive-measure regular product set $\mathcal P$.
Let $r_{\mathcal P}$ be its actual first-return time under $T_R$.
There are finitely many pairs $(n_i,k_i)\in\Z_{>0}\times\Z^2$ such that
\begin{equation}\label{eq:positive-return-witness}
 E_i=\{x\in\mathcal P:r_{\mathcal P}(x)=n_i,
                  S_{n_i,R}\kappa_R(x)=k_i\},\qquad \nu(E_i)>0,
\end{equation}
and the vectors $g_i=(k_{i,1},k_{i,2},n_i)$ generate $\Z^3$ over $\Z$.
In particular, integers $b_{ji}$ can be chosen with
$e_j=\sum_i b_{ji}g_i$ for $1\le j\le3$.
\end{lemma}
\begin{proof}
Let $L$ be the subgroup generated by all positive-measure cells in
\eqref{eq:positive-return-witness}. The cells form a countable partition
of $\mathcal P$ modulo a null set. If $L$ were proper,
Lemma~\ref{lem:integer-subgroup-separation} would give a nontrivial
finite character $\theta=(u,s)\in\T^3$ annihilating every cell record.
For $\phi=u\cdot\kappa_R+s$ this means
$e^{iS_{r_{\mathcal P},R}\phi}=1$ on $\mathcal P$ almost everywhere.

The first-return tower over $\mathcal P$ represents almost every point
of the invertible ergodic collision system uniquely as $T_R^j y$,
where $y\in\mathcal P$ and $0\le j<r_{\mathcal P}(y)$.
Set $q(T_R^j y)=e^{iS_{j,R}\phi(y)}$ on these measurable tower levels.
The preceding return identity verifies the phase equation also at
the top of each level. This is an exact measurable physical phase at
$(u,s,0)\ne0$, contradicting
Theorem~\ref{thm:complete-physical-phase}. Thus $L=\Z^3$.
Writing the three standard basis vectors as integer combinations uses
only finitely many cell records, which gives the asserted family.
The witnesses have positive measure; they are not isolated orbits.
\end{proof}
